% ======================================================================
% 08_introduction.tex: CHAPTER 1, introduction
% WHAT TO WRITE
%   Rationale: why the topic matters and what problem exists.
%   Scientific and Practical Significance: what is new and who benefits.
%   Objectives: the general goal and a list of specific objectives (the list below is a bullet list).
%   Object and Scope: what is studied and what is not.
%   Research Methods: how the work is done.
%   Thesis Structure: one line per chapter. The \ref{chap:...} give the chapter numbers automatically.
%
% HOW TO USE THIS FILE
%   - Replace every ... with your text. A new paragraph starts after one blank line.
%   - The title in \chapter{...} is printed as "CHAPTER n: TITLE" in capitals. Rename it freely.
%   - \section{...} and \subsection{...} are numbered automatically (1.1, 1.1.1). Rename, copy or delete them.
%   - \paragraph{Small heading:} gives a bold line of its own (the colon is required).
%   - The last section ("Template Examples") shows a figure, an equation and a table. Copy a block into your
%     text, edit it, then delete the "Template Examples" section when you no longer need it.
%   - Syntax for everything (lists, figures, tables, equations, citations): docs/LATEX_CHEATSHEET.md
%   - Labels (\label{...}) must be unique across the whole report.
% ======================================================================

\chapter{INTRODUCTION}
\label{chap:introduction}

\section{Rationale}
...

\section{Scientific and Practical Significance}
...

\section{Objectives}
General objective: ...

Specific objectives:
\begin{itemize}
    \item ...
    \item ...
    \item ...
\end{itemize}

\section{Object and Scope}
...

\section{Research Methods}
...

\section{Thesis Structure}
The thesis has the following chapters:
\begin{enumerate}
    \item Chapter~\ref{chap:introduction}: ...
    \item Chapter~\ref{chap:background}: ...
    \item Chapter~\ref{chap:design}: ...
    \item Chapter~\ref{chap:experiments}: ...
    \item Chapter~\ref{chap:results}: ...
    \item Chapter~\ref{chap:conclusion}: ...
\end{enumerate}

% ============================================================
% Template examples: copy these blocks into the relevant subsection.
% Replace placeholder text, captions, labels, and sample data as needed.
% ============================================================
\section{Template Examples}
\label{sec:template-ch1-introduction}

\subsection{Figure and Cross-Reference Example}
...

Figure~\ref{fig:template-ch1-introduction} illustrates ...
\begin{figure}[htbp]
    \centering
    % Replace the framed placeholder with your image, for example:
    % \includegraphics[width=0.75\linewidth]{figures/example.png}
    \fbox{%
        \begin{minipage}[c][4cm][c]{0.75\linewidth}
            \centering
            ...
        \end{minipage}%
    }
    \caption{...}
    \label{fig:template-ch1-introduction}
\end{figure}

\subsection{Equation Example}
...

An example relationship is given in Equation~\eqref{eq:template-ch1-introduction}:
\begin{equation}
    y(t) = a\,x(t) + b,
    \label{eq:template-ch1-introduction}
\end{equation}
where $x(t)$ is the input, $y(t)$ is the output, and $a$ and $b$ are
example coefficients. Replace this relationship and its definitions with ...

\subsection{Table Example}
...

Table~\ref{tab:template-ch1-introduction} presents ...
\begin{table}[htbp]
    \centering
    \caption{...}
    \label{tab:template-ch1-introduction}
    \begin{tabularx}{0.9\linewidth}{lXX}
        \toprule
        \textbf{Item} & \textbf{Description} & \textbf{Value / Unit} \\
        \midrule
        ... & ... & ... \\
        ... & ... & ... \\
        \bottomrule
    \end{tabularx}
\end{table}
